\subsection{重複固有値 $-1$ の位数4枝の有限表示と正常核}
\label{subsec:cubic-fourth-root-negative-repeat-finite}
\begin{theorem}[残る重複 $-1$ の二スペクトルも空]
\label{thm:cubic-fourth-root-negative-repeat-finite-empty}
三固有値必要リストの $(4;2,0,1)$ と $(4;2,0,3)$ はH3を持たない。
この関係パターンのAS4整域の中山は半単純でなければならない。
\end{theorem}
\begin{proof}
まず $\mu=(-1,-1,1,i)$ を半単純とし、$U=\langle x,y\rangle$ とする。
必要関係重複度は $(2,1,2,1)$。
最高tensorの右巡回軌道は17個であり、その代表を
\[
 xxxx,xxxy,xxyy,xxbb,xyxy,xyyy,xybb,xbyb,xbby,
 xbcc,xcbc,xccb,yyyy,yybb,ybcc,ycbc,yccb
\]
とし、係数を $a_0,\ldots,a_{16}$ とする。
重み $1$ の語順 $(xx,xy,yx,yy,bb)$、重み $-1$ の語順
$(xb,yb,bx,by,cc)$、重み $-i$ の語順 $(xc,yc,cx,cy)$ と重み $i$ の $(bc,cb)$ では
\[
 H_0=\begin{pmatrix}
 a_0&a_1&a_1&a_2&a_3\\
 a_1&a_4&a_2&a_5&a_6\\
 a_1&a_2&a_4&a_5&a_8\\
 a_2&a_5&a_5&a_{12}&a_{13}\\
 a_3&a_6&a_8&a_{13}&0
 \end{pmatrix},\quad
 H_2=\begin{pmatrix}
 0&a_7&a_3&a_8&a_9\\
 -a_7&0&a_6&a_{13}&a_{14}\\
 -a_3&-a_6&0&-a_7&-a_{11}\\
 -a_8&-a_{13}&a_7&0&-a_{16}\\
 -a_9&-a_{14}&a_{11}&a_{16}&0
 \end{pmatrix},
\]
\[
 H_{3,1}=\begin{pmatrix}a_{10}&a_{11}\\a_{15}&a_{16}\\
 ia_9&ia_{10}\\ia_{14}&ia_{15}\end{pmatrix}.
\]
前節と同様、$bc,cb$ 関係の両係数非零から第二列は第一列の $q\ne0$ 倍。
$U$ 基底を選んで $a_9=1,a_{14}=0$ と正規化すると
\[
 a_{10}=q,\ a_{11}=q^2,\ a_{15}=a_{16}=0.
\]
$H_2$ のPfaffianは
\[
 -a_{13}a_3+a_6a_8-a_7^2,\quad a_6-a_7q^2,\quad
 a_{13},\quad -a_7+a_8q^2,\quad a_{13}q^2.
\]
従って $a_{13}=0,a_6=q^4a_8,a_7=q^2a_8$。
$H_0$ の $bb$ 列は純 $U^2$ 関係
\[
 r=a_3x^2+a_8q^4xy+a_8yx
\]
を与える。$a_8=0$ なら、この列は純 $x^2$ 又は零である。
前者は整域性に反し、後者は $I_0$ 全体が $U^2$ に含まれて二本の二元二次関係を持つため整域性に反する。
よって $a_8\ne0$。
また $I_0\cap U^2=kr$ である。
$yy$ 列 $(a_2,a_5,a_5,a_{12},0)^t$ もこの直線に入るから
\[
 a_{12}=0,\quad a_2=a_3a_5/a_8,\quad a_5(1-q^4)=0.
\]
$a_5=0$ なら、全二次関係を $(x,b)$ で商にすると
$k\langle y,c\rangle/(c^2)$ となり、前節と同じ指数増大の矛盾が生じる。
従って $a_5\ne0,q^4=1$。
$H_0$ の $xy,yx$ 列の差から $(a_4-a_2)(xy-yx)$ も純関係となるため、
対称な $r$ と独立にならないためには $a_4=a_2$ が必要である。
$y$ に $x$ の倍数を加えれば $a_3=0$ とでき、そのとき $a_2=a_4=0$。
$H_0$ の $(xx,xy,bb)$ 主minorは $-a_0a_8^2$ となり、$a_0=0$。

従って全六関係は、ある $C,d\ne0$ と $A\in k$ に対して
\begin{align*}
 &yx+xy=0,\quad b^2+A x^2+C y^2=0,\quad bx-q^2xb=0,\\
 &by+q^2yb+d c^2=0,\quad cx-iqxc=0,\quad cb+q^{-1}bc=0.
\end{align*}
独立な生成元のscalar変更で $C=d=1$、$A\in\{0,1\}$ に正規化できる。
実際、$x'=r x,y'=s y,b'=t b,c'=u c$ として
$s^2=Ct^2,u^2=dts$、さらに $A\ne0$ なら $r^2=At^2$ と選べばよい。
従って半単純の必要候補は $q\in\{1,-1,i,-i\}$ と $A\in\{0,1\}$ の八表示で覆われる。
この段階では八表示全てのAS4十分性を主張していない。

公開exact検算は八つそれぞれの元の全二次関係から
$B_3,B_4$ の関係行列を作り、商次元 $20,35$ を得る。
$\mu$-不変な三次自由基底六本について $Gv-\mu^{-1}(v)G$ の四次商像を計算すると、
\[
 W_1(B)=\begin{cases}k\,c^2x&q=1,\\0&q=-1,i,-i.\end{cases}
\]
この有限個のexact線型消去は公開コードとcertificateにより再生できる。
$q=1$ の二表示でも全 $G$ は $c\cdot(cx)$ と因子分解するため、整域切断は存在しない。

単冪部分を許しても候補は増えない。
前節のように $H=\mu_u^{1/4}$ で半単純AS4整域 $D=B^H$ へuntwistする。
$H$ は $D$ の自己同型でもあり、$\mu_s$ と可換、$b,c$ を固定する。
一次元の関係 $I_{-i}\subset Uc+cU$ は $x$ の直線を一意に定めるので、
単冪性から $H(x)=x,H(y)=y+t x$。
純 $U^2$ 関係 $xy+yx$ の保存は $2t x^2=0$ を与え、$t=0$。
従って $H=1$、もとの中山も半単純であり、既述の八表示に尽くされる。
正常三次 $G$ のtwist輸送はここでも仮定していない。
$i\mapsto-i$ の場合にも同じ係数恒等式と消去が適用される。
\end{proof}

\begin{corollary}[三固有値の七個の必要スペクトル]
\label{cor:cubic-three-eigenvalue-7-list}
三固有値H3の必要リストは位数3の3個、位数4の4個、計7個、Galoisラベル4軌道である。
代表は $(3;0,1,2),(3;1,0,2),(4;1,0,3),(4;1,2,3)$。
残存候補の存在・許容性及び全H3の分類は未完成である。
\end{corollary}
